% ch16: 原文 GLDraft_D.tex 行 3803–4099；附录 C（3803–3854）与附录 D（3855–4099）
% 只译散文（标题、正文、图题、脚注）；公式、\label、\ref、\eqref、\cite、图文件名保持不变。

\section{典范 oper 截面之间关系的一些例子}\label{somex}
\subsection{$\SL(2,\C)$ 的情形}

基本表示有一组由典范截面（canonical section）构成的基：$s_2 \equiv s$、$Ds$，其中 $D^2 s = t s$。\zy{译者按（依据 \cite{BD} 的一般理论）：这里的“典范截面”指基本截面 $s$ 的对称幂 $s^{\otimes n}$；其各阶导数由 $s$ 与 $Ds$ 的对称化张量积 $s^{\otimes(n-k)}Ds^{\otimes k}$ 给出。在 $\SL(2,\C)$ 情形，一个 oper 由 $C$ 上的一个射影结构连同 $K^{1/2}$（即自旋结构）的一个选择组成；选定自旋结构后，$D^2s=ts$ 就是相应的二阶微分方程。}
$\SL(2,\C)$ 的不可约表示都可以写成基本表示的对称积。
对称幂 $s^{\otimes n}$ 给出 $(n+1)$ 维表示中的典范截面 $s_{n+1}$。$s_{n+1}$ 满足什么微分方程？
在把每个 $D^2 s$ 都换成 $t s$ 之后，对象 $D^r s_{n+1}$（$r\leq n$）可以写成 $s^{\otimes (n-k)} Ds ^{\otimes k}$（$k\leq r$）的对称化形式的线性组合。
到 $n+1$ 阶时不再出现新的对象，于是得到一个微分方程。例如，
\begin{equation}
D^3 s_3= D^2(s \otimes Ds + Ds \otimes s) =2 D(Ds \otimes Ds + 2 t s \otimes s) = 6 t Ds_3 + 4 Dt s_3.
\end{equation}


在把表示相乘时，反对称化可以利用 $s \wedge Ds=1$ 化简。例如，有
\begin{equation}
s_2 \otimes Ds_2 = \frac12 D s_3 + \frac12 s_1
\end{equation}
其中 $s_1=1$，并且我们省略了 $2 \otimes 2 = 3\oplus 1$ 的嵌入映射。我们甚至可以写成一个算子乘积展开（OPE）：
\begin{equation}
s_2(z) \otimes s_2(w) = s_3(w) + \frac12(z-w)Ds_3(w) + \cdots - \frac12 (z-w) s_1(w)+ \cdots
\end{equation}
 
\subsection{$\SL(3,\C)$ 的情形}
基本表示有一组由截面构成的基：$s_3 \equiv s$、$Ds$、$D^2s$，其中 $D^3 s = t_2 Ds +t_3 s$。
外积 $s_{\bar 3} = s_3 \wedge Ds_3$ 是反基本表示中的一个典范截面。确实有
\begin{equation}
D^3s_{\bar 3} = D^2(s_3 \wedge D^2s_3) = D(D s_3 \wedge D^2s_3+t_2 s_3 \wedge D s_3) = t_2 Ds_{\bar 3} + (D t_2-t_3) s_{\bar 3}
\end{equation}

平方 $s_{6,5} = s_3 \otimes s_3$ 是对称平方表示的一个典范截面。
前五个导数 $D^n s_{6,5}$ 是独立的。第六个导数可以用较低阶导数表示，但系数是 $t_2$、$t_3$ 及其导数的有理函数。
更好的做法是定义第二个基元素
\begin{equation}
s_{6,1} \equiv s_3 \otimes D^2 s_3 + D^2 s_3 \otimes s_3 - D s_3 \otimes Ds_3- t_2 s_3 \otimes s_3
\end{equation} 
于是得到如下形式的关系
\begin{align}
&D^5 s_{6,5}- 5 t_2 D^3 s_{6,5}+ (7 t_3 + 4 D t_2)D^2 s_{6,5} + \cdots = (4 t_3 - 2 D t_2) s_{6,1} \cr
&D s_{6,1} = (2 t_3- D t_2) s_{6,5} 
\end{align} 
\zy{译者按（无外部依据）：由下文可见，$s_{6,5}$、$s_{6,1}$、$s_{8,5}$、$s_{8,3}$ 的第二个下标是相应截面所满足的微分方程的阶（亦即其独立导数的个数）：例如 $s_{6,5}$ 满足 5 阶方程，$s_{8,3}$ 满足 3 阶方程。}
我们可以把这些关系以及更多关系概括进一个 OPE
\begin{equation}
s_3(z) \otimes s_3(w) \sim s_{6,5} + \cdots + \frac12 (z-w) s_{\bar 3} +\cdots + \frac14 (z-w)^2 s_{6,1} + \cdots
\end{equation}

注意 $s_3$ 与 $s_{\bar 3}$ 彼此正交，因此 $s_{8,5} = s_3 \otimes s_{\bar 3}$ 是一个典范的伴随截面，而
 $s_{8,3} = s_3 \otimes D s_{\bar 3}-D s_3 \otimes s_{\bar 3}$ 是第二个。它们满足预期形式的微分方程
 $D^5 s_{8,5} = \cdots$ 与 $D^3 s_{8,3}= \cdots$，并且两个方程的右端都同时出现这两个截面的导数。
 
 
 \section{波函数奇异行为的局部模型}\label{localmodel}
 
 \subsection{摇摆除子}\label{wobbly}
 
如第 \ref{wkb} 节所讨论，Hecke 算符或 Hitchin 哈密顿量的本征函数 $\Psi$ 自然地由一张秩 1 拉格朗日膜 $\B_F$ 表示，该膜支撑在 Hitchin 纤维化的一个纤维 $F$ 上。$F$ 必须满足量子修正后的 WKB 条件。
 
为了把 $\Psi$ 表示成 $\M(G,C)$ 上的显式半密度，可以取它与 delta 函数态 $\delta(x,x_0)$ 的内积，得到 $\Psi(x_0)=\la \delta(x,x_0),\Psi\ra$。delta 函数态 $\delta(x,x_0)$ 由一张拉格朗日膜 $\B_{x_0}$ 表示，该膜支撑在余切
丛 $T^*\M(G,C)\to \M(G,C)$ 于 $x_0\in \M(G,C)$ 处的纤维 $L_{x_0}$ 上。尽管内积 $\la \delta(x,x_0),\Psi\ra$ 的计算是一个困难问题，多半不会有相当显式的答案，但它是良定义的；只要 $\M_H(G,C)$ 的 $A$-模型中的膜 $\B_{x_0}$ 良定义并随 $x_0$ 实解析地变化，所得结果就会随 $x_0$ 实解析地变化。而这又只要
$L_{x_0}$
在 $\M_H(G,C)$ 中是闭的便成立。
 
然而已知，这一结论沿一个除子 $\sD\subset \M(G,C)$ 失效，该除子有时被称为摇摆除子（wobbly divisor；参见例如 \cite{DP,PP,PPN,HZ}）。稳定丛 $E\to C$ 称为非常稳定的（very stable），如果任何幂零的希格斯场 $\varphi:E\to E\otimes K$ 实际上都是 0。反之，若存在某个幂零且非零的 $\varphi:E\to E\otimes K$，则称 $E$ 是摇摆的。除子 $\sD$ 参数化摇摆丛。\zy{译者按：原文所引 \cite{PP} 的原题为“The Wobbly Divisors of the Moduli Space of Rank-2 Vector Bundles”，\cite{PPN} 则讨论“Very Stable Bundles and Properness of the Hitchin Map”（Hitchin 映射的紧合性）。“摇摆”与“非常稳定”是互补的：非常稳定丛不允许非零幂零希格斯场，摇摆丛则允许。}
 
膜的支撑集必须是闭的，因此当 $L_{x_0}$ 不是闭集时，它就不能作为膜的支撑集。于是波函数 $\Psi(x_0)$ 可能沿 $\sD$ 变奇异。
 
在本附录中，我们将对 $G=\SU(2)$ 的情形分析沿 $\sD$ 的行为；稍作改动，讨论也适用于 $G=\SO(3)$。对于带抛物结构的丛，亏格 0 情形已有关于 Hecke 算符与 Hitchin 哈密顿量本征函数行为的显式计算 \cite{EFK3}。尽管我们在本文中没有显式纳入抛物结构，在本附录中也不会这样做，但我们的全部考虑都自然推广到带抛物结构的丛。（在规范场论中，抛物结构是通过纳入某一类与余维 2 奇点相伴的面算符来描述的 \cite{GuW2}。）因此，我们将把我们的结果与 \cite{EFK3} 中已得到的结果作比较。
 
以下 $C$ 是亏格为 $g$ 的黎曼曲面，其典范丛为 $K$、切丛为 $T$；$E\to C$ 是秩 2 的全纯向量丛，$\det\,E$ 平凡。
若 $E$ 不是非常稳定的，则存在非零希格斯场 $\varphi:E\to E\otimes K$ 使 $\varphi^2=0$。此时 $E$ 有直线子丛 $\L=\mathrm{ker}\,\varphi$，因此 $E$ 是一个扩张
\be\label{extform} 0\to \L \to  E\to \L^{-1}\to 0.\ee
我们关心非平凡扩张，但希望与直和 $E_0=\L\oplus \L^{-1}$ 作比较。我们主要关心的情形将是：$E_0$ 不稳定，而 $E$ 稳定。若满足下述条件，这一点对一般情形成立：\footnote{次数为 0 的秩 2 丛 $E\to C$ 稳定，当且仅当它的任何秩 1 子丛 $\M\subset E$ 的次数都为负。
若 $E$ 有一个次数为 0 的直线子丛，但没有正次数的直线子丛，则称它是半稳定的。
按定义，正合列 (\ref{extform}) 中的 $E$ 有子丛 $\L$；若 $\L$ 的次数为负，这不会使 $E$ 不稳定（而且出现在这种正合列中的 $E$ 一般都是稳定的）。然而，当 $\L$ 的次数为负时，$\L^{-1}$ 的次数为正，直和 $E_0=\L\oplus \L^{-1}$ 不稳定。}
$$c_1(\L)<0. $$
对一般的 $\L$，$E_0$ 作为 $\SL(2,\C)$ 丛的自同构群是由对角矩阵构成的群，我们把它记作 $P$。

此情形下的幂零希格斯场为 $\varphi_0\in H^0(C,\Hom(\L^{-1},\L\otimes K))=H^0(C, \L^2\otimes K)$。可以预期，在摇摆除子 $\sD$ 的一般点上，幂零希格斯场的空间是 1 维的，因此我们预计会需要\footnote{$H^q(C,\M)$ 的维数记作 $h^q(C,\M)$，或简记作 $h^q(\M)$。} $h^0(\L^2\otimes K)= 1$。
为简单起见，我们将研究 $\sD$ 的一个分支，它具有如下性质：在 $\sD$ 内，$\L$ 可以一般地形变（只固定 $c_1(\L)$）。
（$\sD$ 中存在一些分支，上述性质对它们不成立；这些分支可以用类似方式研究，结果也类似。）
在这一假设下，要求一般情形下 $h^0(\L^2\otimes K)$ 为 1，
就推出一般情形下 $h^1(\L^2\otimes K)$ 为 0，而 $\L$ 的次数由 Riemann–Roch 公式确定：
\be\label{degl} 1=h^0(\L^2\otimes K)-h^1(\L^2\otimes K) = 1-g+ c_1(\L^2\otimes K)=g-1 +2c_1(\L)\ee
于是
\be\label{egf}c_1(\L)=1-\frac{g}{2},\ee
因此，如上所述的构造只对偶数 $g$ 可能。不过，对奇数 $g$，我们可以做非常类似的事情：$c_1(\L)$ 的半整数值
意味着我们应当取规范群为 $\SO(3)$ 而非 $\SU(2)$，并把 $E$ 处处替换为 $\ad(E)$。
由于记号上更简单，我们将继续使用
$\SU(2)$ 规范场论的语言，但下面的评注也适用于另一情形。
对 $g>2$，$c_1(\L)<0$，因此如式 (\ref{extform}) 那样的扩张 $E$ 可以是稳定的，而直和 $E_0$ 不稳定。对 $g=2$，$E$ 可以稳定，而 $E_0$ 半稳定。\zy{译者按（无外部依据）：$g$ 为奇数时 $c_1(\L)=1-g/2$ 是半整数，不可能是一个 $\SU(2)$ 丛的直线子丛的次数；把规范群换成 $\SO(3)=\SU(2)/\Z_2$ 后，$\ad(E)$ 是相应的伴随三维向量丛（伴随表示即其标准三维表示），$\L$ 则只以“分数次”线丛的意义出现。}

条件 $h^0(\L^2\otimes K)=1$ 与 $h^1(\L^{-2})=1$ 互为 Serre 对偶。由于 $H^1(C,\L^{-2})$ 参数化 $E_0$ 的{\it 不}
保持三角结构 (\ref{extform}) 的形变（而不保持这一结构的 $E$ 的形变计数相同），条件
$h^1(\L^{-2})=1$ 确实是我们所需要的，它使得我们所讨论的结构在 $\M(G,C)$ 中以余维 1 出现，从而可以代表沿 $\sD$ 的一般行为。
现在也可以数出 $E_0$ 的{\it 确实}具有式 (\ref{extform}) 三角结构的形变。
这个形变空间的维数是
$h^1(\L^2)=g-1-2c_1(\L) =2g-3$。作为检验，通过形变 $E_0$ 得到的稳定丛族的维数是
$h^1(\ad(E_0)) - h^0(\ad(E_0))-\mathrm{dim}\,P = g + (2g-3)+1 -1 =3g-3$，正是预期值。（我们计入了 $E_0$ 的 $g$ 个对角形变。
在计算 $E_0$ 的形变参数空间时，
我们减去 ${\mathrm{dim}}\,P=1$，以计入一般形变会破坏 $E_0$ 的 $P$-不变性这一事实。）
为下文方便，定义
\be\label{weddo} N = h^1(\L^2)= 2g-3,\ee
即沿除子 $\sD$ 的不稳定丛 $E_0$ 的上三角部分的形变数目。因此，如上所述的构造只给出 $N$ 的正奇数值；
不过，正偶数值可以通过带抛物结构的类似构造，或通过放弃“$\L$ 可以在保持在 $\sD$ 内的同时任意形变”这一假设来得到。



现在我们来考察此情形下希格斯丛模空间的局部行为。
取基使 $E_0=\begin{pmatrix}\L\cr \L^{-1}\end{pmatrix}, $ 则不稳定丛 $E_0$ 可以由一个 $\bar\partial$ 算子显式描述：
\be\label{barpar}\bar D=\bar\partial +\begin{pmatrix} a_0 & 0 \cr 0 & -a_0\end{pmatrix}. \ee
下文中，改变 $\L$ 的“对角”形变并不重要，因为我们已假设任何一般的 $\L$（具有适当的次数）
都能导致我们所考虑的情形。换言之，我们的假设使得 $\L$ 的形变只是让我们沿着 $\sD$ 移动。
出于同样理由，我们对对角希格斯场也不感兴趣。
相反，我们必须考察“上对角”与“下对角”形变，它们使我们离开 $\sD$。
按我们的假设，上对角丛形变的空间
维数为 $N$，下对角丛形变的空间维数为 1。我们取 $H^1(C,\L^2)$ 的一组基 $b_1,\cdots, b_n$，并把
该群的一般元素写成 $\sum t_i b_i$，其中 $t_i$ 为复参数。类似地，取非零的 $c\in H^1(C,\L^{-2})$，并把一般
元素写成 $uc$，其中 $u$ 为复参数。
于是 $\bar D$ 的一般非对角扰动给出
\be\label{arpar}\bar D'=\bar D +\begin{pmatrix} 0 & \sum_{i=1}^n t_i b_i \cr uc & 0  \end{pmatrix}. \ee
类似地，对希格斯场，引入非零元素 $e$（取自 $H^0(C, K\otimes \L^2)$）以及 $H^0(C,K\otimes \L^{-2})$ 的一组基 $f_1,f_2,\cdots,f_N$，
再引入额外的复参数 $s$ 与 $r_1,r_2,\cdots, r_N$，则希格斯场的非对角部分为
\be\label{offed}\varphi_\perp =\begin{pmatrix} 0 & se \cr \sum_{i=1}^N r_i f_i&0 \end{pmatrix}. \ee
对这些数据，需要施加希格斯丛方程 $[\bar D',\varphi_\perp]=0$。由于由 $\varphi_\perp$ 的定义
已有 $[\bar D,\varphi_\perp]=0$，该条件约化为 $\mu=0$，其中
\be\label{kinny} \mu=\sum_{i=1}^N r_i t_i -su. \ee 
（这里通过适当归一化波函数，把某些常数取成了 1。）
为了得到希格斯丛模空间的局部描述，更确切地说，得到它横截于那些无关紧要的对角形变的一个切片，我们还
必须除以 $P\cong \C^*$，即不稳定丛 $E_0$ 的自同构群。$P$ 以如下权作用在这些变量上：上对角
参数 $t_i$ 与 $s$ 的权为 1，下对角参数 $u$ 与 $r_i$ 的权为 $-1$。除以 $P$ 并施加条件
(\ref{kinny}) 的联合操作有一个简单的解释。引入辛形式
\be\label{inny}\Omega = \sum_{i=1}^N \d r_i\, \d t_i +\d s \,\d u, \ee 
便看出 $\mu$ 是 $P$ 作用在这些参数上的矩映射。令 $\mu=0$ 并除以 $P$
的联合操作是一个复辛商，其结果将是一个复辛流形，正如希格斯丛模空间所应有的那样。

然而，我们必须小心“商”的含义。若一个希格斯丛中所有权为 1 的参数 $t_i$ 与 $s$ 都为零，
则它不稳定。\footnote{设 $E$ 是行列式平凡、秩为 2 的丛。若 $E$ 有一个正次数的 $\varphi$-不变子丛，则希格斯丛 $(E,\varphi)$ 不稳定。若 $E$ 有一个次数为零的 $\varphi$-不变子丛，但没有正次数的此类子丛，则称 $(E,\varphi)$ 半稳定。
若 $E$ 有一个正次数子丛，但 $E$ 的任何 $\varphi$-不变子丛次数都为负，则丛 $E$ 不稳定，而希格斯丛 $(E,\varphi)$ 稳定。在目前的例子中，若 $t_i=s=0$，则 $E$ 有正次数的 $\varphi$-不变子丛 $\L^{-1}$。}因此我们希望带上 Mumford 稳定性条件来取商，使得 $t_i$ 与 $s$ 不允许全部为零。于是，在 $P$ 的作用之下，$N+1$ 个变量
$t_i$ 与 $s$ 定义 $\bCP^N$ 中的一个点，而由此构造得到的希格斯丛模空间的局部模型便是
\be\label{hb} M_H= T^*\bCP^N. \ee
更确切地说，这是横截切片，抓住了希格斯丛模空间在 $\sD$ 附近的相关方面。

那么，半稳定丛的模空间 $M$ 的相应局部模型又如何呢？为此我们忘掉希格斯参数 $s$ 与 $r_i$，
只保留丛参数 $t_i$ 与 $u$。若令所有 $t_i$ 为零，便得到不稳定丛。因此 $t_i$ 取为不全为零，而在 $P$ 的作用下，它们定义 $\bCP^{N-1}$ 中的一个元素。
把权为 $-1$ 的变量 $u$ 考虑进来，便看出
$M$ 的局部模型是线丛 $\O(-1)\to \bCP^{N-1}$ 的全空间。这与在一点处爆破 $\C^N$ 相同。
不是非常稳定的丛构成的除子 $\sD$ 在这个模型中由方程 $u=0$ 定义；换言之，它就是例外除子
$\bCP^{N-1}\subset M$。方便的做法是定义 $x_i= u t_i$。于是 $M$ 由 $x_i$ 参数化，而点 $\vec x=0$ 被爆破。
至于与 $x_i$ 典范共轭的变量，可以取 $p_i=r_i/u$，因为
\be\label{wobo}\Omega =\sum_{i=1}^N \d p_i \d x_i. \ee
当然，这些坐标在例外除子附近并不好用。

$N=1$ 的情形在许多方面都是例外。
对 $N=1$，丛 $E_0$ 是半稳定的而非不稳定的，因为 $\L$ 的次数为 0。此外，对 $N=1$，刚才描述的构造中实际上没有爆破。只有当 $N>1$ 时，除子
$\sD\cong \bCP^{N-1}\subset M$ 才是摇摆的。对 $N=1$，它是半稳定的；也就是说，它参数化一族半稳定丛。

现在我们可以解释本附录开头所作的断言：沿摇摆
除子，丛模空间的余切丛的纤维在希格斯丛模空间中不是闭的。在这一分析中，我们
把丛模空间 $\M(G,C)$ 与 $\M_H(G,C)$ 替换为它们的局部模型 $M$ 与 $M_H$。

首先考虑 $M$ 中不在 $\sD$ 内的一个点。$\sD$ 在 $M$ 中的补集由非零的 $N$ 元组 $\vec x=(x_1,x_2,\cdots, x_N)$ 构成。
例如，取由 $x_1\not=0$、$x_2=\cdots= x_N=0$ 定义的点 $p$。由于 $x_1=ut_1$，$x_1\not=0$
意味着 $u$ 与 $t_1$ 都非零。我们可以用 $P$ 的作用固定
$t_1=1$。希格斯场受以下条件约束：
\be\label{higgscon} r_1-su=0, \ee
而对 $r_2,\cdots, r_N$ 没有约束。在利用式 (\ref{higgscon}) 解出 $r_1$（用 $s$ 表示）之后，我们看到 $p$ 处的希格斯场空间
由 $s,r_2, \cdots, r_N$ 参数化。这是预期的 $\C^N$；它是 $M$ 的余切丛在 $p$ 处的纤维，并且在 $\M_H$ 中是闭的。我们将把它记作 $L_p$。现在换而考察
余切丛在点 $p'\in \sD$ 处的纤维。为此，只需取 $u=0$，仍令 $t_1=1$ 以及 $t_i=0,$ $i>1$。式 (\ref{higgscon})
仍然成立，并告诉我们令 $r_1=0$，而对 $u$ 没有约束。因此在这种情况下，余切纤维 $L_{p'}$ 仍由
$s$ 与 $r_2,\cdots, r_N$ 参数化。但 $L_{p'}$ 在 $M_H$ 中不是闭的。为看出这一点，考虑 $\C\subset L_{p'} $ 的一个拷贝，它由 $(t_1,s)=(1,\lambda),\,\lambda\in\C$
与 $r_2=\cdots = r_N=0$ 定义。在 $P$ 的作用下，条件 $(t_1,s)=(1,\lambda)$ 等价于 $(t_1,s)=(1/\lambda,1)$。
显然，极限 $\lambda\to\infty$ 存在。
因此这个 $\C$ 的拷贝在 $M_H$ 中不是闭的；它的闭包包含一个 $(t_1,s)=(0,1)$ 的点，从而把 $\C$ 紧化为 $\bCP^1$。更一般地，$L_{p'}$ 在 $M_H$ 中的闭包
由 $(t_1,s,r_2,\cdots, r_N)$ 模去 $P$ 的作用来参数化，并带有一个约束：$t_1$ 与 $s$ 不允许同时为零。它是 $\bCP^1$ 上的一个 $\C^{n-1}$ 丛。满足
$t_1=0$、$s\not=0$ 的点对应于这样的希格斯丛 $(E_0,\varphi)$：底层丛 $E_0$ 不稳定，但配对 $(E_0,\varphi)$
稳定。这些点包含在 $M_H$ 中，但不包含在 $T^*M$ 中。若令 $t_1=s=0$，便得到不稳定的希格斯丛。

现在让我们回到一般点 $p\in M$（它不包含在 $\sD$ 内），并问当 $p$ 趋近 $\sD$ 时 $L_p$ 会怎样。我们考虑与
前面相同的点 $p$，但既然现在知道允许 $t_1=0$ 很重要，我们就不令 $t_1=1$，而把约束 (\ref{higgscon}) 写成：
\be\label{iggscon} t_1 r_1-su=0. \ee
我们还有点 $p$ 的定义：
\be\label{defp} ut_1= x_1.\ee
$L_p $ 由满足这些方程的 $t_1,u,s$ 与 $r_1,r_2\cdots, r_N$ 参数化，再模去 $P$ 的作用。
为了取 $p\to p'$ 时 $L_p$ 的极限，只需在式 (\ref{defp}) 中令 $x_1=0$。所得簇有两个分支。在一个分支上，
$u=r_1=0$，而 $t_1,s$ 一般非零。这个分支（我们称之为 $L_{p',1}$）就是前面描述过的 $L_{p'}$ 在 $M_H$ 中的闭包；它是 $\bCP^1$ 上的一个 $\C^{N-1}$ 丛。在第二个分支上，$t_1=0$，于是 $s$ 必须非零（为使 $(E,\varphi)$ 稳定），从而式
(\ref{iggscon}) 蕴含 $u=0$。第二个分支是 $\C^N$ 的一个拷贝，由 $r_1,\cdots, r_N$ 参数化。由于 $t_1=u=0$，这第二个分支 $L_{p',2}$
完全包含在 $M_H$ 中 $T^*M$ 的补集内。

对一般的 $p$，存在一张秩 1 的 $A$-膜 $\B_p$ 支撑在 $L_p$ 上；由于 $L_p \cong \C^N$ 单连通，它是唯一的。当 $p\to p'$ 时，
$L_p$ 分裂成两个分支，其中任何一个都可以支撑一张 $A$-膜。由于 $\B_p$ 只在远离摇摆除子处光滑变化，
Hitchin 哈密顿量的本征函数有可能沿该除子变奇异。
我们将在第 \ref{singular} 节研究这个问题。

对 $N=1$，没有摇摆除子，上述构造中也没有爆破。（亏格 2 的模空间确实有摇摆除子，但不具有
我们所假设的那些性质。）$N$ 元组 $(t_1,\cdots, t_N)$ 退化为单个变量
$t$，而模空间 $M$ 在这个局部模型中是由 $x=ut$ 参数化的 $\C$ 的一个拷贝。$x=0$ 处的行为是例外的，但原因
不同。在 $x=0$ 处，丛 $E$ 退化为直和 $E_0=\L\oplus \L^{-1}$，而它在 $N=1$ 时是半稳定的。沿除子 $\sD$ 存在未破缺的 $\U(1)$ 规范对称性，在理解沿
该除子的行为时必须把它考虑进去。

\subsection{波函数的奇异行为}\label{singular}





一般的 Hitchin 哈密顿量是 $\int_C \alpha \Tr \varphi^2$，
其中 $\alpha\in H^1(C,T)$。由于式 (\ref{offed}) 中希格斯场 $\varphi$ 的形式，任何这样的哈密顿量对 $s$ 是 1 次齐次的，同样对 $r_i$ 也是 1 次齐次的。由于线性无关的 Hitchin 哈密顿量的数目
与希格斯丛参数的数目相同，我们可以一般地选取一组 Hitchin 哈密顿量，使得
\be\label{genbas} H_i = s r_i,~~~i=1,\cdots,N.\ee（由于还存在我们忽略掉的对角扰动，这
并不是 Hitchin 哈密顿量的全集。它是在这一近似下依赖于非对角扰动的那一组 Hitchin 哈密顿量。）

用远离原点处好用的典范变量 $p_i, x_i$ 表示，它便是 $H_i=p_i\sum_{k=1}^N x_k p_k$。量子化之后，
它成为
\be\label{qut} H_i=\frac{1}{2}\sum_{k=1}^n \left(\frac{\partial}{\partial x_i} x_k \frac{\partial}{\partial x_k}+\frac{\partial}{\partial x_k} x_k\frac{\partial}{\partial x_i}\right). \ee
选取这样的算符次序，是为了使这些算符形式对称（已知 Hitchin 哈密顿量具有这一性质）。
$H_i$ 作用在如下半密度上：
\be\label{murky} F(x_1,x_2,\cdots, x_N)(\d x_1\d x_2\cdots \d x_N)^{1/2}. \ee
当然，式 (\ref{qut}) 中写下的并不是 Hitchin 哈密顿量的精确公式。
它只是对小 $\vec x$ 有效的渐近公式。考虑一个标度变换，使 $x, p$ 的权分别为 $1,-1$。
式 (\ref{qut}) 中的算符 $H_i$ 以权 $-1$ 标度。这里的主张是：Hitchin 哈密顿量由这些表达式给出，模去
非负权的项。

要用这个模型预测 Hitchin 哈密顿量本征函数在例外除子 $\sD$ 附近的行为，我们需要使用
在 $\sD$ 附近好用的坐标。一个方便的选择是使用 $x_1$ 以及 $y_k=x_k/x_1$（$k=2,3,\cdots,N$）。这些坐标
在 $\sD$ 中一个大的开集附近是好用的。

除了把变量从 $x_1,x_2,x_3,\cdots,x_N$ 换为 $x_1,y_2,y_3,\cdots, y_N$，我们还希望以一种在这些坐标下自然的方式表示半密度。
一个适当的形式是
\be\label{urky} G(x_1,y_2,\cdots, y_N) (\d x_1 \d y_2 \d y_3\cdots \d y_N)^{1/2}. \ee
于是 $F$ 与 $G$ 之间的关系应为
\be\label{rky} F(x_1,x_2,\cdots,x_N)=x_1^{-(N-1)/2} G(x_1,y_2,\cdots, y_N), \ee
从而 $F(x_1,\cdots, x_N)(\d x_1\cdots \d x_N)^{1/2}=G(x_1,y_2,\cdots, y_N) (\d x_1 \d y_2\cdots \d y_N)^{1/2}$。
作用在 $G$ 上的 Hitchin 哈密顿量在新坐标下成为
\begin{align} H_1& = \frac{\partial}{\partial x_1}x_1\frac{\partial}{\partial x_1} -\frac{1}{2}\sum_{k=2}^N\left( y_k \frac{\partial}{\partial y_k}+\frac{\partial}{\partial y_k} y_k\right)
\frac{\partial}{\partial x_1} \cr
H_k& =\frac{\partial^2}{\partial x_1 \partial y_k},~~~~~~~~ k>1. \end{align}
寻找在 $x_1\to 0$ 时表现为 $x_1^\alpha$ 的联合本征函数，可以看出：对偶数 $N$，可能值为 $\alpha=0$ 与 $\alpha=(N-1)/2$，
这与亏格 0 情形抛物丛的计算 \cite{EFK3} 一致。对于奇异解 $x_1^{(N-1)/2}g(x_1,\vec y)$（其中 $g(x_1,\vec y)$
在 $x_1=0$ 处全纯），$g(x_1,\vec y)$ 在 $x_1=0$ 处实际上与 $\vec y$ 无关。对奇数 $N$，相应的陈述是：除全纯
解之外，
还有形如
  $f(x_1,\vec y)+g(x_1,\vec y)x_1^{(N-1)/2}\log x_1$ 的解，其中 $f(x_1,\vec y)$ 与 $g(x_1,\vec y)$ 在 $x_1=0$ 处全纯，并且同样 $g(x_1,\vec y)$
  在 $x_1=0$ 处与 $\vec y$ 无关。
  
  
 
这些公式只抓住了除子 $D$ 附近最奇异的一些项，行为相当非一般。
特别地，本征值方程 $H_s\Psi=\lambda_s\Psi$（$s=1,\cdots, N$）没有预期的 $2^N$ 个线性无关联合解。
若给 Hitchin 哈密顿量加上一般的高阶项而保持其交换性，预期的行为就会恢复。
例如（把 $\partial/\partial x_1$ 简写为 $\partial_1$，把 $\partial/\partial y_k$ 简写为 $\partial'_k$），一个使 Hitchin
哈密顿量的交换性保持到 $\alpha$ 一阶的简单扰动，就是把上面的公式替换为
\begin{align}\label{perturbed} \t H_1& = \partial_1 x_1\partial_1 -\frac{1}{2}\sum_{k=2}^N(y_k \partial'_k+\partial'_k y_k)\partial_1
+2 \sum_{l,m=2}^N \alpha_{klm}\partial'_l y^k \partial'_m \cr
\t H_k&= \partial_1\partial'_k +\sum_{l,m=2}^N \alpha_{klm}\partial'_l\partial'_m,~~~~~k>1, \end{align}  
其中 $\alpha_{klm}=\alpha_{kml}$ 是一般的复系数。奇异行为与前面类似。

\begin{yizhu}[译者按：原文笔误与记号]
本块原文有几处笔误；公式与记号中的笔误在译文中照原样保留，拼写错误则按正确词义译出，不另作径改。
(1) 本块第一式（$D^3s_3=\cdots$）的中间式与末式系数有误。按 $D^2s=ts$ 直接计算，
$D^2(s\otimes Ds+Ds\otimes s)=2D(Ds\otimes Ds+t\,s\otimes s)$，
故 $D^3s_3=4t\,Ds_3+2Dt\,s_3$；原文分别写作 $2t$、$6t$、$4Dt$。
(2) 拼写：“with we will call $P$”应为“which we will call”；“By constrast”应为“By contrast”；“divsor”应为“divisor”；“suppport”应为“support”。
(3) 式 (\ref{arpar}) 前的 $b_1,\cdots, b_n$ 与式中的 $\sum_{i=1}^n$，式 (\ref{qut}) 中的 $\sum_{k=1}^n$，以及 $L_{p'}$ 闭包的纤维 $\C^{n-1}$，按上下文均应为大写 $N$。
(4) “near the divisor $D$”中的 $D$ 应为 $\sD$（摇摆除子）。
(5) 式 (\ref{perturbed}) 第一行中 $\alpha_{klm}\partial'_l y^k\partial'_m$ 的指标 $k$ 未求和，使 $\t H_1$ 依赖于一个未指定的 $k$，疑有笔误；译文照录。
(6) 本块“$h^1(\ad(E_0))-h^0(\ad(E_0))-\dim P=g+(2g-3)+1-1=3g-3$”一式的算术与记号疑不一致：按 $h^0=1$、$\dim P=1$，左边应为 $3g-4$；若只减去 $\dim P$，则为 $3g-3$。译文照录，存疑待核。
\end{yizhu}
